% Options for packages loaded elsewhere
\PassOptionsToPackage{unicode}{hyperref}
\PassOptionsToPackage{hyphens}{url}
%
\documentclass[
]{article}
\usepackage{amsmath,amssymb}
\usepackage{lmodern}
\usepackage{iftex}
\ifPDFTeX
  \usepackage[T1]{fontenc}
  \usepackage[utf8]{inputenc}
  \usepackage{textcomp} % provide euro and other symbols
\else % if luatex or xetex
  \usepackage{unicode-math}
  \defaultfontfeatures{Scale=MatchLowercase}
  \defaultfontfeatures[\rmfamily]{Ligatures=TeX,Scale=1}
\fi
% Use upquote if available, for straight quotes in verbatim environments
\IfFileExists{upquote.sty}{\usepackage{upquote}}{}
\IfFileExists{microtype.sty}{% use microtype if available
  \usepackage[]{microtype}
  \UseMicrotypeSet[protrusion]{basicmath} % disable protrusion for tt fonts
}{}
\makeatletter
\@ifundefined{KOMAClassName}{% if non-KOMA class
  \IfFileExists{parskip.sty}{%
    \usepackage{parskip}
  }{% else
    \setlength{\parindent}{0pt}
    \setlength{\parskip}{6pt plus 2pt minus 1pt}}
}{% if KOMA class
  \KOMAoptions{parskip=half}}
\makeatother
\usepackage{xcolor}
\setlength{\emergencystretch}{3em} % prevent overfull lines
\providecommand{\tightlist}{%
  \setlength{\itemsep}{0pt}\setlength{\parskip}{0pt}}
\setcounter{secnumdepth}{-\maxdimen} % remove section numbering
\ifLuaTeX
  \usepackage{selnolig}  % disable illegal ligatures
\fi
\IfFileExists{bookmark.sty}{\usepackage{bookmark}}{\usepackage{hyperref}}
\IfFileExists{xurl.sty}{\usepackage{xurl}}{} % add URL line breaks if available
\urlstyle{same} % disable monospaced font for URLs
\hypersetup{
  hidelinks,
  pdfcreator={LaTeX via pandoc}}

\author{Dietmar Gerald Schrausser}
\date{09.05.2026}


\begin{document}

\hypertarget{foliumblat-iiir}{%
\section{Foliu{[}m{]}/Blat IIIr}\label{foliumblat-iiir}}

Eine zeitgenössische deutsche Übersetzung wird präsentiert, die auf
einem Vergleich der bestehenden lateinischen, frühneuhochdeutschen
(ENHG) und englischen (Übersetzungs-) Texte basiert, um ein besseres
Verständnis zu ermöglichen. Insbesondere da der ENHG-Text im Vergleich
zum modernen Hochdeutschen schwer verständlich ist, da dieser fundierte
Deutschkenntnisse (als Muttersprache) sowie Kenntnisse verschiedener
Dialekte voraussetzt.

\hypertarget{vom-werk-des-andern-tages}{%
\subsection{Vom Werk des andern Tages}\label{vom-werk-des-andern-tages}}

\begin{quote}
\textbf{S}Ecundo die dixit deus. Fiat firmamentum in medio
aquaru{[}m{]}: et diuidat aquas ab aquis: vocauitq{[}ue{]} firmamentum
celum.
\end{quote}

Lasst uns den Propheten über himmlische Mystetien sprechen hören. Gott
schuf ein \emph{radförmiges}\footnote{`rotundum' und `gescheibelt',
  scheibenförmig.}, bewegliches \emph{Firmament}, in dem sich weitere
\emph{fühlbare}\footnote{`sensibilia' und `empfintliche'.} Dinge
befanden. Aus gefrorenem Wasser erstarrte es zu einem Kristall, in
diesem wurden die Sterne\footnote{`sidera' und `gestirne'.}
angeheftet.\footnote{vgl. das `Speculum maius' (De Beauvais,
  \href{https://nbn-resolving.org/urn:nbn:de:hbz:061:1-20612}{1964},
  \emph{Speculum historiale}, Buch I, Kapitel 1), beschreibt die
  physikalische Beschaffenheit des Firmaments und die \emph{Wasser über
  den Himmeln}, stark beeinflusst von Honorius Augustodunensis' `Imago
  Mundi'
  (\href{https://archive.org/detail/DHonoriiAugustudunensisLibriSeptem/page/n6/mode/1up}{1532}).}
Die Achsen der Sphäre des Himmels, an welcher die Sterne fixiert sind,
sollen aber \emph{zwei} sein. Diese, die (i) nördliche
(\emph{Boreas}\footnote{Griechischer Gott des Nordwinds und
  Personifizierung des Winters.}) und die (ii) südliche
(\emph{Notus}\footnote{\emph{Notus}, der griechische Gott und die
  Personifizierung des Südwindes oder \emph{nothus}, \emph{unehelich},
  \emph{falsch}, \emph{gemischt}, \emph{hybrid}.}), \emph{rollen} mit
solcher Geschwindigkeit von Osten nach Westen, dass dies zum
Zusammenbruch der Welt führte, würden die Planeten diese Bewegung nicht
in entgegengesetzte Richtung abbremsen. Darüberhinaus mäßigte der
Schöpfer der Welt die Natur des Himmels mit Wassern, damit die
Feuersbrunst des Himmels nicht die unteren Elemente in Brand
setze.\footnote{aus dem `Imago Mundi' (Augustudunensis,
  \href{https://archive.org/detail/DHonoriiAugustudunensisLibriSeptem/page/n6/mode/1up}{1532},
  \emph{De Coelo}, Buch I, Kapitel V), einer der populärsten
  Enzyklopädien des Mittelalters.}

{[}\ldots{]}\footnote{`Etlich lerer nit allein auß den unsern sunder
  auch auß den hebreyschen vn{[}d{]} andern haben geglawbt', Etliche
  Lehrer, nicht nur unsere, sondern auch unter den Hebräern und anderen
  glaubten, dass\ldots{}} \emph{Über}\footnote{`Supra', vorher
  vs.~`uber'.} den (I) neun Himmelsorben, das heisst (i) den sieben
Planeten, (ii) der achten Sphäre, der \emph{unirrenden}\footnote{`inerrantem'
  und `die unirrende{[}n{]}'.}, (iii) dem neunten Orbit, durch die
Vernunft, nicht durch die Sinne erfassbar, sowie der erste unter den
\emph{bewegten} Körpern, nimmt man an, dass (II) der zehnte Himmel
unveränderlich, ruhend und still ist, da dieser an keiner Bewegung
teilnimmt.

Nicht nur unser Volk war dieser Ansicht, insbesondere die jüngeren
\textbf{`Strabo'} und \textbf{Beda}, sondern auch viele Hebräer und
ferner einige Philosophen und Mathematiker. Es genügt, wenn zwei, welche
dies belegten Erwähnung finden. Das sind \textbf{Abraham}, der größte
spanische Astrologe und \textbf{Isaak}, der Philosoph. Hier
interpretiert \emph{Isaak} den \emph{zehnten Orbit} als von
\textbf{Hesekiel} durch den \emph{Saphir in Gestalt eines Thrones}
bezeichnet, die Farbe des Saphirs symbolisiert dabei die Brillanz des
Lichts, die Analogie\footnote{`similitudo' und `gleichnus'.} mit einem
Thron symbolisiert die Unbeweglichkeit.\footnote{aus dem `Heptaplus'
  (Mirandola \& Mirandola,
  \href{https://doi.org/10.3931/e-rara-61217}{1601}, Buch II, Kapitel
  1), erörtert die Struktur des Himmels und die Existenz einer zehnten
  Sphäre (des \emph{Empyreums} oder der unbeweglichen \emph{ersten
  Intelligenz}) jenseits der traditionellen neun Sphären.}

⋅

Doch kehren wir zu \emph{Moses} zurück, der \textbf{\emph{inmitten des
Firmaments das Wasser von den Wassern trennte}}, was eine dreifache
Unterteilung der \emph{sublunaren} Körper ergibt. Mehrere befinden sich
(i) oberhalb der mittleren Luftschicht. Im obersten Teil nämlich sind
dessen Elemente und reinstes Feuer, gänzlich der
\emph{Äther}\footnote{`etheris'.}, reine, pure\footnote{`immixta' und
  `unvermischt'.} und wahre\footnote{`legitima' und `richtige'.}
Elemente. Weitere befinden sich (ii) unterhalb des
Mittellbereichs\footnote{`meditulliu{[}m{]}' und `mitteln fürscheinenden
  stat'.} der Luft selbst, also bei uns, wo es kein reines Element (und
auch kein reines fühlbares Element) gibt, sondern ausschliesslich
\emph{Gemenge}\footnote{`mixta' und `gemischt'.}, wegen des
sedimenthaltigen\footnote{`feculenta'.} und dichteren Teils des
weltlichen Körpers, sind.\footnote{vgl. das `Calumniatorem Platonis'
  (Bessarion,
  \href{https://books.google.com/books?id=gPJCAAAAcAAJ}{1503}, Buch IV,
  Kapitel 6), eine neuplatonische Interpretation von Genesis 1,6,
  erklärt die \emph{Teilung der Wasser} durch eine dreifache Aufteilung
  der sublunaren Körper.}

Die (iii) dazwischenliegende Luftschicht, \emph{ebenfalls Firmament}
genannt, woher auch die Vögel stammen, die unter dem Firmament des
Himmels, von diesem im Flug eingeführt werden, ist die Region, in der
die \emph{ätherischen}\footnote{`sublimes'.} Impressionen wie Regen,
Schnee, Blitz, Donner, Kometen oder ähnliches, erscheinen. \emph{Dieses}
Firmament (iii) trennt daher richtigerweise die höheren Elemente (i) von
den niederen (ii), dies nicht nur aufgrund seiner Lage, sondern auch
aufgrund seiner natürlichen Eigenschaft verschiedener \emph{Wasser}.
Über\footnote{`Sup{[}ra{]}' und `darob'.} ihm befinden sich die reinen
Elemente, darunter, \emph{vollständigt} vermengte in elementarer
Einfachheit\footnote{`simplicitate' und `slechtigkeit'.}
abgeschiedene\footnote{`discedunt' und `obgesünderet'.}. So benannte er
dieses Firmament, von Koylon (konkav) ableitend als
\emph{Himmel}\footnote{`celum' (koylon) und `himel'.}, da dieser alle
fühlbaren\footnote{`sensibilia' und `empfintliche'.} und unsichtbaren
Dinge umfasst.\footnote{vgl. das `Epitoma in Almagestum Ptolemaei'
  (Regiomontanus, \href{https://doi.org/10.3931/e-rara-528}{1496},
  \emph{Praefatio}), spiegelt eine mittelalterliche Synthese biblischer
  Kosmologie und aristotelischer Physik wider.}

\hypertarget{references}{%
\subsection{References}\label{references}}

Augustudunensis, H. (1532). \emph{D. Honorii Augustudunensis Presbyteri
Libri Septem}. Basileae.
\url{https://archive.org/detail/DHonoriiAugustudunensisLibriSeptem/page/n6/mode/1up}

Bessarion. (1503). \emph{In Calumniatorem Platonis \ldots: Libri IV}.
Venetiis: Aldus Manutius.
\url{https://books.google.com/books?id=gPJCAAAAcAAJ}

De Beauvais, V. (1964). \emph{Speculum maius}. Graz: Akad. Druck- u.
Verl.-Anst. \url{https://nbn-resolving.org/urn:nbn:de:hbz:061:1-20612}

Mirandola, G., \& Mirandola, G. F. (1601). \emph{Ioannis Pici,
Mirandulae \ldots{} opera quae extant omnia \ldots{}} Basileae: per
Sebastianum Henricpetri. \url{https://doi.org/10.3931/e-rara-61217}

Regiomontanus. (1496). \emph{Epytoma Ioannis de Monte Regio in
almagestum Ptolemei}. Venetiis: Johannes Hamman.
\url{https://doi.org/10.3931/e-rara-528}

\end{document}
